\section{The unconditional cubic bound, after Mirandola and Z\'emor}
The coordinatewise-product theorem can be applied before imposing double centring. Keeping the row and column sums as additional weighted coordinates gives a smaller unconditional dimension bound. The weights have both signs, so a stabilizer argument must control their possible zero-sum blocks.

\begin{lemma}[Weighted moments and the product stabilizer]
\label{lem:weighted-moment-stabilizer}
Let $w_1,\ldots,w_M\in\mathbb C^*$ and let $V\subseteq\mathbb C^M$
be a linear space on which
\[
 \sum_{\nu=1}^M w_\nu v_\nu
 =\sum_{\nu=1}^M w_\nu v_\nu^2
 =\sum_{\nu=1}^M w_\nu v_\nu^3=0.
\]
Products of vectors are coordinatewise, and products of subspaces
denote their linear spans. Put
\[
 U=\mathbb C\mathbf1+V,\qquad
 W=U^2,\qquad H=\operatorname{St}(W).
\]
If $H\cap V=0$, then
\[
 3\dim V+\dim H\le M.
\]
In particular $3\dim V\le M-1$.
\end{lemma}
\begin{proof}
Use the nondegenerate bilinear form
$\langle x,y\rangle_w=\sum_\nu w_\nu x_\nu y_\nu$.
Polarization of the three identities gives
$W\subseteq V^{\perp_w}$, hence
\[
 \dim W\le M-\dim V.
\]
The stabilizer $H$ is a unital subalgebra. Since $\mathbf1\in W$,
one has $H\subseteq W$ and $HW=W$.
Set $K=HU$. Then $K$ contains $H+V$, and
\[
 \dim K\ge\dim H+\dim V,\qquad K^2=H^2U^2=W.
\]
The coordinate-product theorem
\citep[Theorem~3.3, p.~7]{MZ} gives
\[
 \dim W=\dim K^2
 \ge2\dim K-\dim\operatorname{St}(K^2)
 =2\dim K-\dim H
 \ge2\dim V+\dim H.
\]
Combining the inequalities proves the assertion. Finally
$\dim H\ge1$ because $\mathbf1\in H$.
\end{proof}

\begin{theorem}[The unconditional cubic bound]
\label{thm:weighted-bound}
For every $n\ge3$, every complex level space of order $n$ satisfies
\[
 3\dim L\le n^2+2n-1.
\]
In particular
$\operatorname{LS}(5)\le11$,
$\operatorname{LS}(6)\le15$, and
$\operatorname{LS}(7)\le20$.
\end{theorem}
\begin{proof}
For $A=(a_{ij})\in L$, put
$r_i=\sum_j a_{ij}$ and $c_j=\sum_i a_{ij}$.
The identity $S_1(A)=0$ gives
$\sum_i r_i=\sum_j c_j=0$.
Define
\[
 b_{ij}=a_{ij}-\frac{r_i+c_j}{2},\qquad
 \alpha=-\frac{n-2}{2n},\qquad
 \lambda=\frac{n^3}{(n-2)^2}.
\]
Map $A$ injectively to
\[
 \Phi(A)=(b_{ij};\,\alpha r_i;\,\alpha c_j)
 \in\mathbb C^M,\qquad M=n^2+2n,
\]
and give the first $n^2$ coordinates weight $1$ and the last $2n$
coordinates weight $-\lambda$.
The map is injective because its last coordinates recover the
marginal sums, after which
$a_{ij}=b_{ij}+(r_i+c_j)/2$.

The quadratic and cubic expansions on $S_1=0$ are
\begin{align*}
 2S_2(A)
 &=\sum_{ij}b_{ij}^2
       -\frac n4\left(\sum_i r_i^2+\sum_j c_j^2\right),\\
 \frac32S_3(A)
 &=\sum_{ij}b_{ij}^3
       +\frac{n-2}{8}\left(\sum_i r_i^3+\sum_j c_j^3\right).
\end{align*}
Here is the cubic calculation explicitly. Put $A_3=\sum_{ij}a_{ij}^3$, $R_3=\sum_i r_i^3+\sum_jc_j^3$, $B=\sum_{ij}a_{ij}r_ic_j$ and $C=\sum_{ij}a_{ij}^2(r_i+c_j)$. Applying the distinct-row and distinct-column partition expansion on three labels gives
\[
 \frac32S_3=A_3-\frac32C+\frac32B+\frac12R_3.
\]
Expanding the cube of $b_{ij}$ gives
\[
 \sum_{ij}b_{ij}^3=A_3-\frac32C+\frac32B+\frac{6-n}{8}R_3.
\]
Indeed $\sum_{ij}a_{ij}(r_i+c_j)^2=R_3+2B$ and $\sum_{ij}(r_i+c_j)^3=nR_3$, since the marginal sums are zero. Their difference proves the displayed cubic identity. Expanding the square gives $\sum b_{ij}^2=\sum a_{ij}^2+(n-4)(\sum r_i^2+\sum c_j^2)/4$, which proves the quadratic identity using \eqref{eq:quad}.

Since $-\lambda\alpha^2=-n/4$ and
$-\lambda\alpha^3=(n-2)/8$, the weighted quadratic and cubic moments
vanish on $V=\Phi(L)$.
Its weighted linear moment vanishes as well: the sum of the
$b$ coordinates and the sum of each marginal coordinate group
are zero. We will also use the unweighted linear constraint
\begin{equation}
 \sum_{\nu=n^2+1}^{M}v_\nu=0\qquad(v\in V).
 \label{eq:weighted-extra-sum}
\end{equation}

Take $U,W,H$ as in Lemma~\ref{lem:weighted-moment-stabilizer}.
It remains to prove $H\cap V=0$.
A unital subalgebra of $\mathbb C^M$ has a basis of disjoint
coordinate projectors, corresponding to a partition into blocks
$B$ \citep[Lemma~2.7, p.~5]{MZ}.
On this basis the restriction of the weighted bilinear form to $H$
is diagonal, with entries
\[
 \omega(B)=\sum_{\nu\in B}w_\nu.
\]
Since $H\subseteq W\subseteq V^{\perp_w}$, every element of
$H\cap V$ belongs to the radical of this restricted form.
It is therefore a linear combination of the indicators of
zero-weight blocks.

If a zero-weight block contains $a$ entry coordinates and $b$
marginal coordinates, then
\[
 a(n-2)^2=bn^3,\qquad a,b>0.
\]
Writing $g=\gcd(n^3,(n-2)^2)$, its counts are positive multiples of
\[
 (a_0,b_0)=
 \left(\frac{n^3}{g},\frac{(n-2)^2}{g}\right).
\]
Here $g=1$ for odd $n$, $g=4$ for $n\equiv0\pmod4$, and $g=8$
for $n\equiv2\pmod4$. For the even cases write $n=2m$ and use $g=4\gcd(2,(m-1)^2)$; for odd $n$, the integers $n$ and $n-2$ are coprime.
The constraint $a_0\le n^2$ rules out such a block except at
$n=4$ or $6$.
At those two orders the minimal counts are respectively
$(16,1)$ and $(27,2)$.
There can therefore be at most one zero-weight block in the
partition.

For every other order the restricted form on $H$ is nondegenerate,
so $H\cap V=0$.
At order four or six its radical is either zero or the line spanned
by the indicator $\mathbf1_B$ of that one block.
But $\mathbf1_B$ has a positive number $b$ of marginal coordinates,
so no nonzero scalar multiple of it satisfies
\eqref{eq:weighted-extra-sum}.
Thus $H\cap V=0$ at these orders too.

Lemma~\ref{lem:weighted-moment-stabilizer} now gives the stronger
intermediate inequality
\[
 3\dim L+\dim H\le n^2+2n,
\]
and $\dim H\ge1$ proves the theorem.
\end{proof}


The conclusion is still quadratic in $n$ and uses no layer beyond $S_3$. At order six it gives the unconditional interval $9\le\LS(6)\le15$. The bound ten in Theorem~\ref{thm:rowcycle} requires a symmetry hypothesis, and no reduction of arbitrary spaces to that class is asserted.
